\chapter{Simulation Plan}
\label{chap:sims}

\section{Philosophy}

The framework's strength is that it should be simulable from first principles.
Rare large vacuum fluctuations, Einstein--Cartan bounces, and thermodynamic
equilibrium are all computable problems. The path from ``interesting framework''
to ``constrained by data'' runs through simulation.

The simulation work is tiered. Each tier is a tractable project of definite
scope; higher tiers build on lower ones. A small group could complete Tiers 1--2
in a year; Tiers 3--4 are multi-year programs. Code lives in \texttt{sims/} and
outputs in \texttt{sims/output/}.

\section{Tier 1: Single-bounce Einstein--Cartan dynamics}

\paragraph{Goal.} Numerically evolve a spherical collapse of fermionic matter
through the bounce. Output: complete characterization of the bounce as a function
of input parameters.

\paragraph{Setup.}
\begin{itemize}
  \item 1D spherical symmetry (radial coordinate $r$, time $t$).
  \item Initial conditions: a uniform ball of matter at moderate density,
    slightly overdense to drive gravitational collapse.
  \item Matter modeled as a Weyssenhoff fluid: a perfect fluid with spin density
    (semiclassical treatment of fermion spin in the continuous limit).
  \item Evolution: Einstein--Cartan field equations with torsion sourced by spin
    density.
  \item Tracking: density, pressure, spin density, metric components, and torsion
    components as functions of $r$ and $t$.
\end{itemize}

\paragraph{Numerical method.} A finite-difference scheme on a radial grid
($\sim 10^3$ to $10^4$ points), with adaptive mesh refinement near the bounce
(where gradients become extreme), implicit time stepping for stability near peak
density, and conservation checks (total mass--energy, total spin, total particle
number).

\paragraph{Expected output.} The bounce happens at $\rho \sim \rhoc \sim
10^{50}\,\mathrm{kg\,m^{-3}}$ (verified); the bounce timescale in proper time;
the bounce duration and profile shape; the maximum density reached as a function
of initial conditions; the equation of state through the bounce; and the energy
balance (kinetic energy in expansion vs.\ energy contained in spin/torsion
fields).

\paragraph{Tools.} Python with NumPy/SciPy for prototyping (optionally Julia for
performance), standard ODE solvers (RK4, implicit methods), custom grid handling
and AMR, and plotting with matplotlib.

\paragraph{Effort estimate.} 2--3 months for a competent graduate student;
$\sim 500$--$1500$ lines of code; validation against existing single-bounce
papers (Brandt, Magueijo, Pop\l{}awski~\citep{poplawski2012,poplawski2016}).

\paragraph{Status of existing work.} Several authors have done pieces of this. No
one has produced a clean, well-documented reference implementation. \emph{A
valuable side-output of this tier is producing such a reference implementation.}

\section{Tier 2: Perturbation evolution through the bounce}

\paragraph{Goal.} Take the output of Tier 1 (a clean bounce profile) and evolve
cosmological perturbations through it. Output: the primordial power spectrum from
the bounce, comparable to CMB observations.

\paragraph{Setup.} Start with the Tier 1 spherical bounce background; add linear
perturbations to the metric and matter fields; evolve them using linearized
Einstein--Cartan equations; track the power spectrum of scalar, vector, and
tensor modes through the bounce; and compute predictions for CMB observables (the
scalar spectral index $n_s$, the tensor-to-scalar ratio $r$, the running of the
spectral index, and non-Gaussianities).

\paragraph{Numerical method.} Mode-by-mode evolution in Fourier space;
long-wavelength modes treated as classical fields and short-wavelength modes
quantum-mechanically (mode functions); Bunch--Davies-like initial conditions at
large negative time (deep in collapse); modes tracked through the bounce and into
post-bounce expansion.

\paragraph{Expected output.} The primordial power-spectrum amplitude and shape,
the spectral tilt, and the tensor-to-scalar ratio, compared to Planck CMB
observations ($n_s = 0.965 \pm 0.004$, $r < 0.06$ upper
limit)~\citep{planck2018} and to $\Lambda$CDM $+$ inflation predictions.

\paragraph{Decisive test.} If the bounce produces a near-scale-invariant spectrum
with $n_s \approx 0.96$ and $r$ within observational bounds, \textbf{the
framework is observationally viable.} If the spectrum is dramatically different,
the framework fails this test. This is the central calculational challenge for
\SCT. Existing work on bounce cosmologies (loop quantum cosmology, ekpyrotic
models) provides templates for similar calculations~\citep{ashtekar2006}.

\paragraph{Effort estimate.} 6--12 months for someone with
cosmological-perturbation-theory background; $\sim 1000$--$3000$ lines of code
building on Tier 1.

\section{Tier 3: Chirality dynamics through the bounce}

\paragraph{Goal.} Quantify the flip $+$ filter mechanism. Predict the chirality
bias of a baby universe given the parent's properties.

\paragraph{Setup.} Use the Tier 1 bounce background; include fermion content
explicitly (not just a Weyssenhoff fluid); track baryon number, lepton number,
and chiral currents through the bounce; include electroweak physics at
temperatures above the EW scale (sphaleron processes); compute the filter
efficiency for chirality; and compute net baryon-number production/preservation.

\paragraph{Decisive tests.} Does the framework predict
$\eta = n_B/n_\gamma \sim 6 \times 10^{-10}$? Does it predict a
dark-matter-to-baryon ratio $\sim 5$? Are these robust to parameter choices, or
do they require fine-tuning?

\paragraph{Effort estimate.} 12--24 months; significant theoretical work in
addition to numerical; coupling Einstein--Cartan to the Standard Model is
non-trivial.

\section{Tier 4: Supraverse population statistics}

\paragraph{Goal.} Simulate the full supraverse: many vacuum fluctuations
producing OG bounces, each seeding lineages of descendant universes. Compute
equilibrium distributions of universe parameters; check that observer-supporting
universes cluster near our observed parameters.

\paragraph{Setup.} Treat the supraverse as a stochastic process: vacuum
fluctuations produce OG bounces at random locations; each OG bounce seeds a
lineage; each universe evolves cosmologically and produces descendants through
internal black-hole formation; track populations across many generations; and
compute distributions of universe mass, age, baryon asymmetry, and chirality
lineage depth.

\paragraph{Decisive tests.} An equilibrium distribution exists and is computable;
observer-supporting universes cluster around some peak; and our observed
parameters sit near the peak (validating the fine-tuning resolution).

\paragraph{Effort estimate.} A multi-year project, several people; requires
Tiers 1--3 as input. Output: the most comprehensive cosmological prediction the
framework can make.

\section{Tier 5: Direct observation comparisons}

\paragraph{Goal.} Use the simulations to make specific predictions for
observational programs (Planck, DESI, Euclid, LiteBIRD, JWST, etc.) and compare
directly: a precise CMB spectrum shape including specific deviations from
$\Lambda$CDM; large-scale-structure clustering with dark matter as parent
influence; the Hubble parameter at different redshifts; galaxy-rotation asymmetry
magnitude from parent rotation; and the dark-energy evolution trajectory.

\paragraph{Decisive tests.} The framework is genuinely tested against existing
and forthcoming observational data. Each prediction either matches or does not.

\paragraph{Effort estimate.} Each comparison is its own project; a multi-decade
research program at full scope.

\section{Minimum viable project}

If we want one project that demonstrates the framework's viability:
\textbf{Tier 2 (perturbation through bounce) with a focus on the CMB power
spectrum.} This is decisive---either the bounce reproduces something close to the
observed CMB spectrum or it does not. The result is publishable in either case.
The effort is bounded at $\sim 12$ months, and the required expertise (numerical
relativity $+$ cosmological perturbation theory) is widely available. This should
be the first priority for any simulation effort.

\section{Visualization parallel track}

Alongside the physics simulations, the visualization work
(Chapter~\ref{chap:viz}) can proceed independently. It does not require Tier 2
outputs to start---it can begin with simpler models of supraverse structure and
refine as physics simulations produce better inputs. The visualization track is
more accessible and might be the right place for someone with
computational/graphics interests to start.

\section{Tooling}

\paragraph{Languages.} Python (prototyping, plotting, analysis; default for
Tier 1); Julia (performance-sensitive numerics; optional for Tiers 1--2,
advisable for Tiers 3--4); C++ with the Einstein Toolkit (if existing numerical
relativity infrastructure is leveraged---heavy but well-tested).

\paragraph{Libraries.} NumPy, SciPy, matplotlib (Python); DifferentialEquations.jl
(Julia); Einstein Toolkit, GRChombo (numerical relativity); CAMB or CLASS (CMB
power spectrum, for comparison to standard predictions).

\paragraph{Hardware.} Tiers 1--2: workstation or single-node cluster
($\sim$ hours--days). Tier 3: small cluster ($\sim$ days--weeks). Tier 4:
significant cluster resources ($\sim$ weeks--months). Tier 5: depends on the
specific calculations.

\paragraph{Storage.} Simulations should produce reproducible outputs in
standardized formats (HDF5 preferred); visualization data should be portable
across plotting backends; code in git, with versioning of major releases.

\section{Reproducibility}

All simulation code should be open source from the start, documented well enough
that others can reproduce results, validated against analytic limits where they
exist, and tested for numerical convergence (refine grid, check stability). The
point is not to ``prove'' \SCT through simulations---the point is to make the
framework testable. Reproducible code lets others check, extend, and challenge
results.
